\documentclass[11pt]{article}

\usepackage[a4paper,margin=2.2cm]{geometry}
\usepackage[T1]{fontenc}
\usepackage[utf8]{inputenc}
\usepackage{lmodern}
\usepackage{amsmath,amssymb,mathtools}
\usepackage{microtype}
\usepackage{enumitem}
\usepackage{titlesec}
\usepackage[hidelinks]{hyperref}

\titleformat{\section}{\Large\bfseries}{\thesection}{0.75em}{}
\titleformat{\subsection}{\large\bfseries}{\thesubsection}{0.75em}{}
\titlespacing*{\section}{0pt}{1.4ex plus .2ex minus .2ex}{0.9ex}
\titlespacing*{\subsection}{0pt}{1.0ex plus .2ex minus .2ex}{0.5ex}

\newcommand{\R}{\mathbb{R}}
\newcommand{\diag}{\operatorname{diag}}
\newcommand{\rank}{\operatorname{rank}}

\begin{document}

\begin{center}
    {\LARGE\bfseries LFR Latent Dimension Reduction via SVD}\\[0.5em]
\end{center}

\vspace{0.7em}

\noindent
The LPV-LFR realization derived in \texttt{LFR-derivation-supervisor.tex} uses
$\Delta(Y) = Y I_6$, meaning six latent channels. This document derives a
reduced realization with fewer channels, following the SVD-based procedure of
R.~Tóth. Two methods are presented. Method~1 applies SVD directly to the full
coupling matrix. Method~2 first reduces $D_{zw}$ before applying SVD. Section~2
is a placeholder to be filled in after Section~1 is verified.

\section{Method 1: SVD on the Full Coupling Matrix}

\subsection{Setup}

The starting point is the LPV-LFR realization from
\texttt{LFR-derivation-supervisor.tex}, with constant matrix $G$ and scheduling
block $\Delta(Y) = Y I_6$. The loop and state equations are
\begin{align}
    z(t) &= C_z\, x(t) + D_{zw}\, w(t) + D_{zu}\, u(t),
    \label{eq:loopz} \\
    w(t) &= Y\, I_6\, z(t),
    \label{eq:loopw} \\
    \dot{x}(t) &= A_x\, x(t) + B_w\, w(t) + B_u\, u(t),
    \label{eq:state}
\end{align}
with the constant matrices (reproduced from the derivation for reference)
\begin{equation}
    C_z =
    \begin{bmatrix}
        -M_0^{-1}K & -M_0^{-1}C \\
        0 & 0
    \end{bmatrix},
    \quad
    D_{zw} =
    \begin{bmatrix}
        -M_0^{-1}M_1 & -M_0^{-1}M_2 \\
        I_3 & 0
    \end{bmatrix},
    \quad
    D_{zu} =
    \begin{bmatrix}
        M_0^{-1} \\
        0
    \end{bmatrix},
    \label{eq:CDmats}
\end{equation}
\begin{equation}
    B_w =
    \begin{bmatrix}
        0 & 0 \\
        -M_0^{-1}M_1 & -M_0^{-1}M_2
    \end{bmatrix},
    \quad
    B_u =
    \begin{bmatrix}
        0 \\ M_0^{-1}
    \end{bmatrix}.
    \label{eq:Bmats}
\end{equation}
The goal is to find a reduced latent dimension $n \leq 6$ such that
$\Delta(Y) = Y I_n$ suffices, with a correspondingly smaller set of loop
signals $\tilde{z}(t), \tilde{w}(t) \in \R^n$.

\subsection{Constructing the Stacked Matrix}

Substituting \eqref{eq:loopz} into \eqref{eq:loopw} gives
\begin{equation}
    w(t) = Y\bigl[C_z\, x(t) + D_{zw}\, w(t) + D_{zu}\, u(t)\bigr]
         = Y\,[C_z,\; D_{zw},\; D_{zu}]
           \begin{bmatrix} x(t) \\ w(t) \\ u(t) \end{bmatrix}.
    \label{eq:wexpanded}
\end{equation}
The state equation \eqref{eq:state} contains $B_w w(t)$. Using
\eqref{eq:wexpanded},
\begin{equation}
    B_w\, w(t)
    = Y\, B_w\,[C_z,\; D_{zw},\; D_{zu}]
      \begin{bmatrix} x(t) \\ w(t) \\ u(t) \end{bmatrix}.
    \label{eq:Bwexpanded}
\end{equation}
Both \eqref{eq:wexpanded} and \eqref{eq:Bwexpanded} depend on
$[x(t);\, w(t);\, u(t)]$ through the same right factor $[C_z, D_{zw}, D_{zu}]$,
scaled by $Y$. Stacking them yields
\begin{equation}
    Y \underbrace{
        \begin{bmatrix}
            B_w \\[2pt] I_6
        \end{bmatrix}
        [C_z,\; D_{zw},\; D_{zu}]
    }_{S\;\in\;\R^{12 \times 15}}
    \begin{bmatrix} x(t) \\ w(t) \\ u(t) \end{bmatrix}
    =
    \begin{bmatrix} B_w\, w(t) \\ w(t) \end{bmatrix}.
    \label{eq:stacked}
\end{equation}
The matrix $S$ encodes all dependence of $[B_w w;\, w]$ on $[x;\, w;\, u]$
through the latent variable. Its rank determines the minimal number of latent
channels required.

\subsection{SVD and Change of Basis}

Since $[B_w;\, I_6]$ has $I_6$ in its lower half, it has full column rank $6$.
Therefore
\begin{equation}
    \rank(S)
    = \rank\!\Bigl(\bigl[B_w;\, I_6\bigr]\,[C_z,\; D_{zw},\; D_{zu}]\Bigr)
    = \rank\!\bigl([C_z,\; D_{zw},\; D_{zu}]\bigr).
    \label{eq:rankS}
\end{equation}
Compute the thin SVD of $S$:
\begin{equation}
    S = U\,\Sigma\,V^\top,
    \qquad
    \Sigma = \diag(\sigma_1, \dots, \sigma_n),
    \quad \sigma_1 \geq \cdots \geq \sigma_n > 0,
    \label{eq:svd}
\end{equation}
where $n = \rank(S) \leq 6$. Split the left factor as
\begin{equation}
    L := U\,\Sigma^{1/2}
    =
    \begin{bmatrix} L_x \\ L_z \end{bmatrix},
    \qquad L_x \in \R^{6\times n},\quad L_z \in \R^{6\times n},
    \label{eq:Lsplit}
\end{equation}
and define
\begin{equation}
    R := \Sigma^{1/2}\,V^\top
    = \bigl[R_x,\; R_w,\; R_u\bigr],
    \qquad R_x \in \R^{n\times 6},\quad R_w \in \R^{n\times 6},\quad R_u \in \R^{n\times 3}.
    \label{eq:Rsplit}
\end{equation}
Note that $LR = U\Sigma V^\top = S$ by construction. Define the reduced latent
variable
\begin{equation}
    \tilde{z}(t) := R\,\begin{bmatrix} x(t) \\ w(t) \\ u(t) \end{bmatrix}
    = R_x\, x(t) + R_w\, w(t) + R_u\, u(t) \;\in \R^n.
    \label{eq:ztilde}
\end{equation}
From $S = LR$ and \eqref{eq:stacked}, the original loop signals satisfy
\begin{equation}
    B_w\, w(t) = L_x\, \tilde{w}(t),
    \qquad
    w(t) = L_z\, \tilde{w}(t),
    \label{eq:wfromwtilde}
\end{equation}
where the reduced scheduling relation is
\begin{equation}
    \tilde{w}(t) = Y\, I_n\, \tilde{z}(t).
    \label{eq:wtilde}
\end{equation}
Substituting \eqref{eq:wfromwtilde} into \eqref{eq:ztilde},
\begin{equation}
    \tilde{z}(t)
    = R_x\, x(t) + R_w\, L_z\, \tilde{w}(t) + R_u\, u(t),
    \label{eq:ztildeloop}
\end{equation}
and substituting \eqref{eq:wfromwtilde} into \eqref{eq:state},
\begin{equation}
    \dot{x}(t) = A_x\, x(t) + L_x\, \tilde{w}(t) + B_u\, u(t).
    \label{eq:statereduced}
\end{equation}
The reduced constant matrix $\tilde{G}$ therefore has entries
\begin{equation}
    \tilde{B}_w = L_x,\quad
    \tilde{C}_z = R_x,\quad
    \tilde{D}_{zw} = R_w\, L_z,\quad
    \tilde{D}_{zu} = R_u,
    \label{eq:Gtilde}
\end{equation}
with $A_x$, $B_u$, $C_y$, $D_{yw}$, $D_{yu}$ unchanged. The scheduling block
becomes $\Delta(Y) = Y I_n$ with $n \leq 6$.

\subsection{Application to the Dual-Gantry Model: Why Method 1 Gives No Reduction}

Method 1 reduces the latent dimension if and only if $\rank([C_z, D_{zw}, D_{zu}]) < 6$.
We now show that for the dual-gantry realization this rank equals $6$, so
Method 1 cannot reduce $\Delta(Y) = Y I_6$.

Write $[C_z, D_{zw}, D_{zu}] \in \R^{6 \times 15}$ by blocks of columns,
grouping state columns $(1$--$6)$, loop columns $(7$--$12)$, and input columns
$(13$--$15)$:
\begin{equation}
    [C_z,\; D_{zw},\; D_{zu}]
    =
    \left[
    \begin{array}{cc|cc|c}
        -M_0^{-1}K & -M_0^{-1}C &
        -M_0^{-1}M_1 & -M_0^{-1}M_2 &
        M_0^{-1} \\[4pt]
        0 & 0 &
        I_3 & 0 &
        0
    \end{array}
    \right].
    \label{eq:CzDzwDzu}
\end{equation}
The top three rows (rows $1$--$3$) and bottom three rows (rows $4$--$6$) are
linearly independent by the following argument.

\paragraph{Step 1: bottom rows have rank 3.}
Restrict rows $4$--$6$ to the loop columns $7$--$9$: the submatrix is $I_3$,
which has rank $3$. Hence rows $4$--$6$ are linearly independent among themselves.

\paragraph{Step 2: bottom rows are zero in the input columns.}
Rows $4$--$6$ restricted to columns $13$--$15$ are the zero matrix.

\paragraph{Step 3: top rows are nonsingular in the input columns.}
Rows $1$--$3$ restricted to columns $13$--$15$ equal $M_0^{-1}$. Since $M_0$
is the inertia matrix evaluated at $Y = 0$ and is positive definite (physical
mass matrix), $M_0^{-1}$ is invertible. Therefore rows $1$--$3$ are linearly
independent among themselves.

\paragraph{Step 4: top and bottom rows are mutually independent.}
Suppose $\alpha^\top(\text{rows } 1\text{--}3) + \beta^\top(\text{rows }
4\text{--}6) = 0$ for some $\alpha, \beta \in \R^3$. Restricting to columns
$13$--$15$:
\begin{equation}
    \alpha^\top M_0^{-1} + \beta^\top \cdot 0 = 0
    \quad\Longrightarrow\quad
    \alpha^\top M_0^{-1} = 0
    \quad\Longrightarrow\quad
    \alpha = 0,
    \label{eq:alpha0}
\end{equation}
where the last implication uses invertibility of $M_0^{-1}$. With $\alpha = 0$,
the equation reduces to $\beta^\top(\text{rows } 4\text{--}6) = 0$, and since
rows $4$--$6$ are linearly independent (Step~1), $\beta = 0$.

\paragraph{Conclusion.}
All six rows are linearly independent, so
\begin{equation}
    \rank\!\bigl([C_z,\; D_{zw},\; D_{zu}]\bigr) = 6.
    \label{eq:rank6}
\end{equation}
From \eqref{eq:rankS}, $\rank(S) = 6$ as well. The SVD of $S$ has six non-zero
singular values, and Method~1 yields $n = 6$: the latent dimension is
unchanged. Method~1 therefore provides no reduction for the dual-gantry
realization.

The reason is structural: $D_{zu} = [M_0^{-1};\, 0]$ introduces a full-rank
$3 \times 3$ block in the input columns of the top rows, while the bottom rows
are zero there. This is a direct consequence of the input $u(t)$ entering the
acceleration equation through $M_0^{-1}$, an invertible matrix. No change of
parameters removes this structure as long as $M_0$ remains invertible.

\section{Method 2: SVD after Pre-Compression of $D_{zw}$}

Method~1 fails because $\rank([C_z, D_{zw}, D_{zu}]) = 6$. The cause is the
$D_{zu} = [M_0^{-1};\, 0]$ block, which is unaffected by any change to $D_{zw}$.
Method~2 first reduces $D_{zw}$ to expose its true rank before applying the SVD
procedure, yielding a loop of dimension $n \leq 4$.

\subsection{Block Structure of $D_{zw}$}

The coupling matrices from the gantry equations are
\begin{equation}
    M_1 =
    \begin{bmatrix}
        0 & -m_h & 0 \\ -m_h & 0 & 0 \\ 0 & 0 & 0
    \end{bmatrix},
    \qquad
    M_2 =
    \begin{bmatrix}
        0 & 0 & 0 \\ 0 & m_h & 0 \\ 0 & 0 & 0
    \end{bmatrix}.
    \label{eq:M1M2}
\end{equation}
From \eqref{eq:CDmats}, $D_{zw}$ decomposes into two row blocks:
\begin{equation}
    D_{zw} =
    \begin{bmatrix}
        D_{\mathrm{top}} \\ D_{\mathrm{bot}}
    \end{bmatrix},
    \qquad
    D_{\mathrm{top}} := -M_0^{-1}[M_1,\, M_2] \in \R^{3\times 6},
    \qquad
    D_{\mathrm{bot}} := [I_3,\; 0] \in \R^{3\times 6}.
    \label{eq:DzwBlocks}
\end{equation}

\subsection{Rank of $D_{zw}$}

Let $(M_j)_k$ denote the $k$-th column of $M_j$. Columns~3, 4, 6 of
$D_{\mathrm{top}}$ are zero because $(M_1)_3 = (M_2)_1 = (M_2)_3 = 0$.
The non-zero columns follow from $(M_1)_1 = [0,\,-m_h,\,0]^\top$,
$(M_1)_2 = [-m_h,\,0,\,0]^\top$, $(M_2)_2 = [0,\,m_h,\,0]^\top$:
\begin{equation}
    (D_{\mathrm{top}})_1 = m_h\,(M_0^{-1})_2,
    \qquad
    (D_{\mathrm{top}})_2 = m_h\,(M_0^{-1})_1,
    \qquad
    (D_{\mathrm{top}})_5 = -m_h\,(M_0^{-1})_2.
    \label{eq:Dtopcols}
\end{equation}
Columns~4 and~6 of $D_{\mathrm{bot}}$ are zero as well, so $\rank(D_{zw}) \leq 4$.
To show that columns~1, 2, 3, 5 are linearly independent, suppose
\begin{equation}
    \alpha_1 c_1 + \alpha_2 c_2 + \alpha_3 c_3 + \alpha_5 c_5 = 0,
    \label{eq:lindep2}
\end{equation}
where $c_k$ denotes column~$k$ of $D_{zw}$. Restricting \eqref{eq:lindep2}
to the bottom block $D_{\mathrm{bot}}$:
\begin{equation}
    \alpha_1 e_1 + \alpha_2 e_2 + \alpha_3 e_3 = 0
    \quad\Longrightarrow\quad
    \alpha_1 = \alpha_2 = \alpha_3 = 0.
\end{equation}
Restricting to the top block with $\alpha_1 = \alpha_2 = \alpha_3 = 0$:
\begin{equation}
    -\alpha_5\, m_h\, (M_0^{-1})_2 = 0.
    \label{eq:alpha5}
\end{equation}
Since $m_h > 0$ and $(M_0^{-1})_2 \neq 0$ (invertibility of $M_0$),
$\alpha_5 = 0$. Therefore
\begin{equation}
    \rank(D_{zw}) = 4.
    \label{eq:rankDzw}
\end{equation}

\subsection{Structural Identity $B_w = E\, D_{zw}$}

From \eqref{eq:Bmats} and \eqref{eq:DzwBlocks}: the bottom block of $B_w$
equals $D_{\mathrm{top}}$, while the top block of $B_w$ is zero. Defining
\begin{equation}
    E :=
    \begin{bmatrix}
        0 & 0 \\ I_3 & 0
    \end{bmatrix}
    \in \R^{6\times 6},
    \qquad
    B_w = E\, D_{zw}.
    \label{eq:BwEDzw}
\end{equation}

\subsection{Thin SVD of $D_{zw}$ and Loop Pre-Compression}

Since $\rank(D_{zw}) = 4$, compute the thin SVD
\begin{equation}
    D_{zw} = \hat{U}\,\hat{\Sigma}\,\hat{V}^\top,
    \qquad
    \hat{U} \in \R^{6\times 4},\quad
    \hat{\Sigma} \in \R^{4\times 4},\quad
    \hat{V}^\top \in \R^{4\times 6},
    \label{eq:DzwSVD}
\end{equation}
and define
\begin{equation}
    S_2 := \hat{U}\,\hat{\Sigma} \in \R^{6\times 4},
    \qquad
    P := \hat{V}^\top \in \R^{4\times 6},
    \label{eq:S2P}
\end{equation}
so that $D_{zw} = S_2\, P$. Define compressed loop variables
\begin{equation}
    \hat{z}(t) := P\, z(t) \in \R^4,
    \qquad
    \hat{w}(t) := P\, w(t) \in \R^4.
    \label{eq:zwhat}
\end{equation}
The scheduling relation $w = Y I_6 z$ projects to
\begin{equation}
    \hat{w}(t) = Y\, I_4\, \hat{z}(t),
    \label{eq:whatnew}
\end{equation}
so the compressed loop retains the same structure with $\Delta(Y) = Y I_4$.
The $D_{zw}$ coupling term in \eqref{eq:loopz} becomes
\begin{equation}
    D_{zw}\, w(t) = S_2\, P\, w(t) = S_2\, \hat{w}(t).
    \label{eq:DzwS2}
\end{equation}

\subsection{Modified Stacked Matrix}

Substituting \eqref{eq:DzwS2} into \eqref{eq:loopz} and multiplying by $P$:
\begin{equation}
    \hat{z}(t) = P C_z\, x(t) + P S_2\, \hat{w}(t) + P D_{zu}\, u(t).
    \label{eq:zhatloop}
\end{equation}
Following the same substitution as Method~1 (replacing $D_{zw}$ by $S_2$ and
$I_6$ by $P$), the modified stacked matrix is
\begin{equation}
    S_{\mathrm{new}} =
    \begin{bmatrix}
        B_w \\[2pt] P
    \end{bmatrix}
    \underbrace{[C_z,\; S_2,\; D_{zu}]}_{6\times 13}
    \;\in\; \R^{10\times 13}.
    \label{eq:Snew}
\end{equation}

\subsection{Rank Bound $n \leq 4$}

From \eqref{eq:BwEDzw} and $D_{zw} = S_2\,P$:
\begin{equation}
    B_w = E\, S_2\, P.
    \label{eq:BwES2P}
\end{equation}
Every row of $B_w$ is a linear combination of rows of $P$, so
\begin{equation}
    \rank\!\begin{bmatrix} B_w \\ P \end{bmatrix} = \rank(P) \leq 4.
    \label{eq:rankBwP}
\end{equation}
It remains to show $\rank([C_z, S_2, D_{zu}]) = 6$. Split $S_2$ conformably
with $D_{zw}$:
\begin{equation}
    S_2 =
    \begin{bmatrix}
        S_{2,\mathrm{top}} \\ S_{2,\mathrm{bot}}
    \end{bmatrix},
    \qquad S_{2,\mathrm{top}},\, S_{2,\mathrm{bot}} \in \R^{3\times 4}.
    \label{eq:S2split}
\end{equation}
From $D_{zw} = S_2\, P$ and \eqref{eq:DzwBlocks}, the bottom block gives
\begin{equation}
    S_{2,\mathrm{bot}}\, P = D_{\mathrm{bot}} = [I_3,\; 0].
    \label{eq:S2botP}
\end{equation}
Since $S_{2,\mathrm{bot}} \in \R^{3\times 4}$ and $\rank(S_{2,\mathrm{bot}}\,P) = 3$:
\begin{equation}
    \rank(S_{2,\mathrm{bot}}) = 3.
    \label{eq:rankS2bot}
\end{equation}
Partition $[C_z, S_2, D_{zu}]$ into top and bottom row blocks conformably
with $D_{zw}$.

\paragraph{Step 1: bottom block has rank~3.}
The bottom block restricted to the $S_2$ columns is $S_{2,\mathrm{bot}}$,
which has rank~3 by \eqref{eq:rankS2bot}.

\paragraph{Step 2: bottom block is zero in the input columns.}
$D_{zu} = [M_0^{-1};\, 0]$ has zero bottom block.

\paragraph{Step 3: top block is nonsingular in the input columns.}
The top block of $D_{zu}$ is $M_0^{-1}$, which is invertible.

\paragraph{Step 4: top and bottom blocks are mutually independent.}
Suppose $\alpha^\top(\text{top block}) + \beta^\top(\text{bottom block}) = 0$
for $\alpha, \beta \in \R^3$. Restricting to the input columns:
\begin{equation}
    \alpha^\top M_0^{-1} + \beta^\top \cdot 0 = 0
    \quad\Longrightarrow\quad
    \alpha^\top M_0^{-1} = 0
    \quad\Longrightarrow\quad
    \alpha = 0.
\end{equation}
Then $\beta^\top(\text{bottom block}) = 0$; by Step~1, $\beta = 0$.

\paragraph{Conclusion.}
All six rows of $[C_z, S_2, D_{zu}]$ are linearly independent, so
$\rank([C_z, S_2, D_{zu}]) = 6$.

Since $[C_z, S_2, D_{zu}]$ has full row rank, right-multiplication by it
preserves rank:
\begin{equation}
    \rank(S_{\mathrm{new}})
    = \rank\!\begin{bmatrix} B_w \\ P \end{bmatrix}
    = \rank(P) \leq 4.
    \label{eq:rankSnew}
\end{equation}
The number of non-zero singular values of $S_{\mathrm{new}}$ is therefore $n \leq 4$.

\subsection{SVD of $S_{\mathrm{new}}$ and Reduced Matrices}

Compute the thin SVD $S_{\mathrm{new}} = U_{\mathrm{new}}\,\Sigma_{\mathrm{new}}\,V_{\mathrm{new}}^\top$
with $n = \rank(S_{\mathrm{new}}) \leq 4$ non-zero singular values. Split
\begin{equation}
    L_{\mathrm{new}} := U_{\mathrm{new}}\,\Sigma_{\mathrm{new}}^{1/2}
    =
    \begin{bmatrix} L_x \\ L_z \end{bmatrix},
    \qquad
    R_{\mathrm{new}} := \Sigma_{\mathrm{new}}^{1/2}\,V_{\mathrm{new}}^\top
    = \bigl[R_x,\; R_{\hat{w}},\; R_u\bigr],
    \label{eq:LRnew}
\end{equation}
where $L_x \in \R^{6\times n}$, $L_z \in \R^{4\times n}$,
$R_x \in \R^{n\times 6}$, $R_{\hat{w}} \in \R^{n\times 4}$, $R_u \in \R^{n\times 3}$.
Define the final reduced latent variable $\tilde{z}(t) := R_{\mathrm{new}}[x;\,\hat{w};\,u] \in \R^n$.
The reduced constant matrix $\tilde{G}$ has entries
\begin{equation}
    \tilde{B}_w = L_x,\quad
    \tilde{C}_z = R_x,\quad
    \tilde{D}_{zw} = R_{\hat{w}}\, L_z,\quad
    \tilde{D}_{zu} = R_u,
    \label{eq:Gtilde2}
\end{equation}
with $A_x$, $B_u$, $C_y$, $D_{yw}$, $D_{yu}$ unchanged, and the final
scheduling block $\Delta(Y) = Y I_n$ with $n \leq 4$.

The computation of $\tilde{G}$ requires the numerical SVD in
\eqref{eq:DzwSVD} and \eqref{eq:LRnew}. The bound $n \leq 4$ is algebraic
and holds for all physically valid parameter values ($m_h > 0$, $M_0 \succ 0$).

\end{document}
